\documentclass[12pt,a4paper]{article}

% -------------------------------------------------
% Basic packages
% -------------------------------------------------
\usepackage[a4paper,margin=1in]{geometry}
\usepackage{fontspec}
\usepackage{ctex}
\usepackage{amsmath,amssymb}
\usepackage{graphicx}
\usepackage{booktabs}
\usepackage{array}
\usepackage{longtable}
\usepackage{caption}
\usepackage{float}
\usepackage{setspace}
\usepackage{fancyhdr}
\usepackage{hyperref}
\usepackage{xcolor}
\usepackage{enumitem}
\usepackage{listings}

% -------------------------------------------------
% Table of contents setup
% -------------------------------------------------
\setcounter{tocdepth}{2}

% -------------------------------------------------
% Font setup
% -------------------------------------------------
\setmainfont{Noto Serif}
\setmonofont{DejaVu Sans Mono}

% -------------------------------------------------
% Line spacing and paragraph formatting
% -------------------------------------------------
\onehalfspacing
\setlength{\parindent}{2em}
\setlength{\parskip}{0.4em}

% -------------------------------------------------
% Header / footer
% -------------------------------------------------
\pagestyle{fancy}
\fancyhf{}
\fancyfoot[C]{\thepage}
\renewcommand{\headrulewidth}{0pt}

% -------------------------------------------------
% Hyperref setup
% -------------------------------------------------
\hypersetup{
    colorlinks=true,
    linkcolor=blue,
    urlcolor=blue,
    citecolor=blue
}

% -------------------------------------------------
% Code listing setup
% -------------------------------------------------
\lstset{
    basicstyle=\ttfamily\small,
    breaklines=true,
    frame=single,
    columns=fullflexible,
    showstringspaces=false
}

% -------------------------------------------------
% User-defined information
% -------------------------------------------------
\newcommand{\HomeworkTitle}{Computational Physics Ex\_3}
\newcommand{\StudentID}{2024141230044}
\newcommand{\StudentName}{范奕}

\begin{document}

% =================================================
% Title block
% =================================================
\begin{center}
    {\LARGE \textbf{\HomeworkTitle}}\\[1.2em]
    {\large Student ID: \StudentID}\\[0.5em]
    {\large Name: \StudentName}
\end{center}

\vspace{0.5em}
\noindent\rule{\textwidth}{0.8pt}
\vspace{0.5em}

\section*{Problem 3 \& Problem 4}

This document records the C programs and the actual execution results for the following two tasks:
\begin{itemize}
    \item Problem 3: check the system endianness of the computer.
    \item Problem 4: check the machine epsilon of \texttt{float}, \texttt{double}, and extended precision (here tested with \texttt{long double}).
\end{itemize}

The source files used in this experiment are:
\begin{itemize}
    \item \texttt{3.c}: solution to Problem 3.
    \item \texttt{4.c}: solution to Problem 4 without using \texttt{float.h}.
    \item \texttt{4f.c}: solution to Problem 4 using constants from \texttt{float.h}.
\end{itemize}

\newpage
\tableofcontents
\newpage

% =================================================
% Part 1
% =================================================
\section{Source Code}

\subsection{Problem 3: Check Endianness}

\begin{lstlisting}[language=C]
#include <stdio.h>

int main() {
    unsigned int x = 1;
    char *p = (char *)&x;

    if (*p == 1) {
        printf("当前系统是小端序（Little Endian）\n");
    } else {
        printf("当前系统是大端序（Big Endian）\n");
    }

    return 0;
}
\end{lstlisting}

\subsection{Problem 4: Compute Machine Epsilon Without \texttt{float.h}}

\begin{lstlisting}[language=C]
#include <stdio.h>

int main() {
    float feps = 1.0f;
    double deps = 1.0;
    long double leps = 1.0L;

    while ((1.0f + feps / 2.0f) != 1.0f) {
        feps /= 2.0f;
    }

    while ((1.0 + deps / 2.0) != 1.0) {
        deps /= 2.0;
    }

    while ((1.0L + leps / 2.0L) != 1.0L) {
        leps /= 2.0L;
    }

    printf("Machine epsilon for float       = %.20e\n", feps);
    printf("Machine epsilon for double      = %.20e\n", deps);
    printf("Machine epsilon for long double = %.20Le\n", leps);

    return 0;
}
\end{lstlisting}

\subsection{Problem 4: Read Machine Epsilon From \texttt{float.h}}

\begin{lstlisting}[language=C]
#include <stdio.h>
#include <float.h>

int main() {
    printf("FLT_EPSILON  = %.20e\n", FLT_EPSILON);
    printf("DBL_EPSILON  = %.20e\n", DBL_EPSILON);
    printf("LDBL_EPSILON = %.20Le\n", LDBL_EPSILON);
    return 0;
}
\end{lstlisting}

% =================================================
% Part 3
% =================================================
\section{Numerical Method}

\subsection*{Problem 3: Principle of Endianness Detection}

The program stores the integer value \texttt{1} in memory and then reads the first byte at its address.

If the first byte is \texttt{1}, then the least significant byte is stored at the lowest address, so the system is \textbf{little endian}. Otherwise, the system is \textbf{big endian}.

In other words, the test is based on the memory layout of an integer.

\subsection*{Problem 4: Definition of Machine Epsilon}

For a floating-point type, machine epsilon is the smallest positive number \(\varepsilon\) such that
\[
1 + \varepsilon \neq 1.
\]

In the program \texttt{4.c}, the value of \(\varepsilon\) is found by starting from \texttt{1.0} and repeatedly dividing by \texttt{2}. When halving once more no longer changes the value of \texttt{1}, the current value is taken as the machine epsilon.

\subsection*{Using \texttt{float.h}}

In the program \texttt{4f.c}, the standard header \texttt{float.h} is used. It directly provides the implementation-defined constants
\texttt{FLT\_EPSILON}, \texttt{DBL\_EPSILON}, and \texttt{LDBL\_EPSILON}.

This provides a convenient way to verify the numerical results obtained in \texttt{4.c}.

% =================================================
% Part 4
% =================================================
\section{Program Output}

Below is the actual output obtained on the computer.

\subsection*{Problem 3 Output}

\begin{lstlisting}
当前系统是小端序（Little Endian）
\end{lstlisting}

Therefore, the current machine uses \textbf{little-endian} byte order.

\subsection*{Problem 4 Output From \texttt{4.c}}

\begin{lstlisting}
Machine epsilon for float       = 1.19209289550781250000e-07
Machine epsilon for double      = 2.22044604925031308085e-16
Machine epsilon for long double = 1.08420217248550443401e-19
\end{lstlisting}

\subsection*{Problem 4 Output From \texttt{4f.c}}

\begin{lstlisting}
FLT_EPSILON  = 1.19209289550781250000e-07
DBL_EPSILON  = 2.22044604925031308085e-16
LDBL_EPSILON = 1.08420217248550443401e-19
\end{lstlisting}

The results from \texttt{4.c} and \texttt{4f.c} are identical, which confirms that the iterative computation is correct.

% =================================================
% Part 5
% =================================================
\section{Discussion}

\subsection*{Discussion for Problem 3}

The result shows that this computer is little endian. This means that for multi-byte data types, the low-order byte is stored at the low memory address.

This is the most common byte order on modern x86 and x86\_64 systems.

\subsection*{Discussion for Problem 4}

From the output, the machine epsilon values are:
\[
\varepsilon_{\mathrm{float}} = 1.1920928955078125 \times 10^{-7},
\]
\[
\varepsilon_{\mathrm{double}} = 2.22044604925031308085 \times 10^{-16},
\]
\[
\varepsilon_{\mathrm{long\ double}} = 1.08420217248550443401 \times 10^{-19}.
\]

These values show that \texttt{double} has much higher precision than \texttt{float}, and \texttt{long double} has even finer precision on this platform.

It should also be noted that the assignment mentions \emph{quadruple precision}, but in standard C the type used here is \texttt{long double}. On different platforms, \texttt{long double} may or may not correspond to true quadruple precision. Therefore, the third result should be interpreted as the machine epsilon of the platform's \texttt{long double} implementation.

\subsection*{Overall Conclusion}

The experiment verifies two basic properties of the current computer system:
\begin{itemize}
    \item The machine is \textbf{little endian}.
    \item The machine epsilon values for \texttt{float}, \texttt{double}, and \texttt{long double} were successfully obtained.
    \item The direct iterative method and the constants in \texttt{float.h} give the same results, confirming the correctness of the program.
\end{itemize}

\vfill
\begin{center}
    \textit{End of Report}
\end{center}

\end{document}
